% ECONOMICS_PROBLEM: {"id":"C73-49","title":"Dynamic games with incomplete information","jel_code":"C73","source_id":"OP-0004"}
% FIRST_POSTED: 2026-10-10
% ENTRY_KIND: research_attempt
% ATTEMPT_STATUS: solved
\documentclass[11pt,a4paper]{article}
\usepackage[margin=25mm]{geometry}
\usepackage{amsmath,amssymb,amsthm}
\usepackage[hidelinks]{hyperref}
\hypersetup{pdftitle={Dynamic games with incomplete information},pdfauthor={Ji Ho Bae}}
\setlength{\parindent}{0pt}
\setlength{\parskip}{4pt}
\setlength{\emergencystretch}{2em}
\allowdisplaybreaks[2]
\numberwithin{equation}{section}
\newtheorem{theorem}{Theorem}
\newtheorem{proposition}[theorem]{Proposition}
\newtheorem{lemma}[theorem]{Lemma}
\newtheorem{corollary}[theorem]{Corollary}
\theoremstyle{definition}
\newtheorem{example}[theorem]{Example}
\newtheorem{remark}[theorem]{Remark}
\newcommand{\R}{\mathbb R}
\newcommand{\Q}{\mathbb Q}
\newcommand{\E}{\mathbb E}
\newcommand{\Prob}{\mathbb P}
\newcommand{\cE}{\mathcal E}
\newcommand{\ind}{\mathbf1}
\DeclareMathOperator{\SE}{SE}
\DeclareMathOperator{\diag}{diag}
\DeclareMathOperator{\diam}{diam}
\DeclareMathOperator{\Reg}{Reg}
\DeclareMathOperator{\excess}{excess}
\DeclareMathOperator{\TV}{TV}
\title{Dynamic Games with Incomplete Information\\[4pt]
 \large Sharp history-based identification and certified inference}
\author{Ji Ho Bae\thanks{Prepared with GPT Codex assistance; the exact
serving revision is unavailable. All mathematical arguments are hand
proofs. No Lean formalization, external peer review, or novelty claim
is asserted.}}
\date{10 October 2026}
\begin{document}
\maketitle
\begin{center}
\small\textbf{C73-49 \quad JEL C73 \quad Source OP-0004}\\
Complete hand proof with internal mathematical and source-fidelity review.
\end{center}
\begin{abstract}
For arbitrary finite known perfect-recall trees, we state restrictions
under which incomplete-information games admit informative, sharp,
and computationally certified inference. A calibrated measurement
channel on complete terminal histories recovers the latent law.
Positive history mass then identifies behavior and beliefs and reduces
the entire compatible parameter set to an explicit affine slice of
a joint parameter polytope. A rational projection and linear program
give finite-sample whole-set confidence, conditioned mean-square
parameter bounds, expected diameter bounds, and an actual feasible
Bayes assessment with a certificate for every continuation deviation.
General policy ranges retain all sequential equilibria at a shared
structural parameter and have exact real-algebraic endpoint certificates.
For two-player zero-sum utility values, sequence-form linear programs
and a coupling bound give polynomial explicit-tree computation and
stable confidence intervals. A tie example proves that the same
contraction claim cannot hold for arbitrary policy outcomes, at any
fixed confidence level. Private histories are retained throughout.
The algorithms are proved; a full QE or LP execution is not claimed.
\end{abstract}

\section{Declared class, observation design, and scope}\label{sec:model}
The canonical entry is an editorial research-family question. It asks
which restrictions make identified and counterfactual sets informative,
and requests verified computation, approximation, sampling confidence,
and an explicit policy-selection convention. The results below give a
complete response for the following declared class. They do not assert
identification under every unrepresented primitive family.

The input is an arbitrary finite known extensive-form tree with perfect
recall. It may be the complete history tree of a finite-horizon game
with persistent private states, correlated signals, and arbitrary
observation histories. Strategic information sets and their common
nonempty action sets are known. No public Markov reduction is made.
Let \(Z\) be the terminal-history set, \(L=|Z|\), \(N_{\rm tr}\) the
number of nodes, and \(H\) the maximum number of edges on a terminal path.
Each \(z\) records the \emph{complete} latent history, including private
signals and Nature events, not merely public outcomes.

The structural parameter is
\(\gamma=(\alpha,\beta)\in\Gamma\subset\R^d\), where \(\Gamma\) is a
nonempty compact rational polytope and \(d\) is the ambient coordinate
count. Joint restrictions on \(\alpha,\beta\) are retained. Empty
parameter domains are detected and reported separately.
The vector of all included Nature-edge probabilities is
\begin{equation}\label{eq:primitives}
 c(\alpha)=c^0+B\alpha,\qquad
 u_{i,z}(\beta)=u^0_{i,z}+\varphi_{i,z}\cdot\beta .
\end{equation}
All displayed coefficients are known and rational. On \(\Gamma\),
Nature rows are proper probability vectors and every included edge
is uniformly positive. Structurally absent edges are omitted; strategic
off-path actions are retained. Repeated transition and signal kernels
are imposed by repeated rows or primitive equalities, not fitted as
unrelated probabilities.

The known discount factor of the underlying dynamic game is denoted
\(\delta_0\), to distinguish it from the unknown payoff vector \(\beta\).
Discounted stage-utility features are summed along a terminal history
to obtain \eqref{eq:primitives}. Cardinal units, known utility entries,
and any normalization are fixed as primitives. Additive or scale
directions not removed by those restrictions remain unidentified.
Let \(d_\beta\) be the dimension of \(\beta\), and suppose
\begin{equation}\label{eq:payoffbounds}
 \|\beta\|_1\le B_1,\quad
 U_0=\max_{i,z}|u^0_{i,z}|,\quad
 F=\max_{i,z,j}|\varphi_{i,z,j}|,\quad U=U_0+FB_1 .
\end{equation}
Empty maxima are zero. Thus \(|u_{i,z}(\beta)|\le U\) on \(\Gamma\).

After history \(z\), the researcher observes one of \(K\) cells through
a known column-stochastic channel \(M(o\mid z)\). This calibrated
post-play measurement is conditionally determined by \(z\); it changes
neither players' information nor their payoffs during play.
Calibration aligns latent labels with their transition and payoff
features. If \(q\) is the terminal law, the observed law is
\begin{equation}\label{eq:channel}p=Mq.\end{equation}
The generic exact method allows rank-deficient \(M\) and zero components
of \(q\). The informative branch assumes \(M\) has full column rank and
the recovered population \(q\) is strictly positive. Quantitative
statements additionally use a declared known rational \(\lambda>0\)
with \(q_z\ge\lambda\), necessarily \(L\lambda\le1\).

Independent markets share one baseline assessment, with independent
Nature and behavioral randomization across markets. An observed
selection stratum can be analyzed separately. An unobserved mixture
over different assessments is a different sampling model.
The sample size is \(n\) throughout this paper.

For each declared policy \(\ell\), another known finite perfect-recall
tree satisfies the same compact, affine, positive-Nature assumptions
at the \emph{same} \(\gamma\). A declared scalar outcome has the form
\begin{equation}\label{eq:policyoutcome}
 W_\ell(\gamma,\sigma)
   =\sum_{z\in Z_\ell}q^\ell_z(\gamma,\sigma)w_{\ell,z}(\gamma),
\end{equation}
where \(w_{\ell,z}\) are known rational polynomials. This includes
probabilities, counts, and calibrated affine utility or welfare
statistics. The attainment assertions concern these continuous
outcomes. Vector targets use the joint image, not a Cartesian product
of coordinate ranges.

Exact algorithms require rational or effectively represented algebraic
numerical inputs, including the channel and numerical population law.
Population statements also make sense for fixed known real inputs;
an arbitrary unencoded real is not finite computational input.
The polynomial-bit claims use rational encodings. The complete tree
and channel may be exponentially large in a succinct horizon description.

The prior catalogue attempt \cite{Prior} already proves the
fixed-positive-Nature exact equilibrium reduction, compactness, sharp
policy projections, endpoint bisection, and whole-set sampling coverage.
We reproduce those arguments and attribute them. The additional content
is the general informative inverse and affine fiber, its quantitative
statistical and actual-assessment certificates, and qualified stable
policy-value inference. Classical sequential equilibrium
\cite{KW}, real quantifier elimination \cite{Basu}, sequence form
\cite{KMS}, and rational linear programming are named inputs;
no priority claim is made.

\section{Exact sequential equilibrium, existence, and compactness}
\label{sec:se}
Fix a baseline or policy tree. Let \(\mathcal B\) be its product of
behavioral simplexes, \(\mathcal B^+\) its strictly positive part,
and \(\mathcal U\) the product of belief simplexes over the nodes of
each information set. Write
\[
 \begin{gathered}
 R_h(\gamma,b)=\prod_{\text{edges on the prefix of }h}
                     \text{their Nature or behavioral probabilities},\\
 R_I(\gamma,b)=\sum_{h\in I}R_h(\gamma,b).
 \end{gathered}
\]
For \(b\in\mathcal B^+\), every \(R_I\) is positive. Its Bayes graph is
\[
 b\in\mathcal B^+,\quad \nu\in\mathcal U,\quad
 R_I(\gamma,b)>0,\quad
 \nu_{I,h}R_I(\gamma,b)=R_h(\gamma,b)\quad(h\in I).
\]
The following formula defines consistency:
\begin{equation}\label{eq:consistency}
 \forall r>0\ \exists(b,\nu)\text{ in this Bayes graph}:
       \|b-\sigma\|_2^2+\|\nu-\mu\|_2^2<r^2 .
\end{equation}
Here \(\gamma\) and all game primitives are held fixed.
One common pair \((b,\nu)\) serves every information set.
The witnesses need not be equilibria or rational.

For player \(i\), information set \(I\), and complete pure continuation
plan \(a_i\), let \(V_{i,I}(\gamma,a_i,\sigma_{-i};\mu)\) be the
terminal-utility sum over suffix paths starting at \(h\in I\), weighted
by \(\mu_{I,h}\). The corresponding prescribed-continuation value is
\(V_{i,I}(\gamma,\sigma;\mu)\). These are finite polynomials.
Require the legitimate-assessment simplex conditions and
\begin{equation}\label{eq:sr}
 V_{i,I}(\gamma,\sigma;\mu)\ge
 V_{i,I}(\gamma,a_i,\sigma_{-i};\mu)
       \quad\text{for every }i,I,a_i .
\end{equation}
Their conjunction with \eqref{eq:consistency} is denoted
\(\SE(\gamma,\sigma,\mu)\).

\begin{lemma}[Exact fixed-game predicate]\label{lem:exactse}
The predicate \(\SE\) describes exactly the sequential-equilibrium
assessments, including consistent off-path beliefs.
\end{lemma}
\begin{proof}
A converging consistency sequence supplies a witness for every \(r>0\).
Conversely, choose witnesses for \(r=1/m\). They give a single sequence
of completely mixed profiles with jointly convergent Bayes beliefs
in the same game.

Perfect recall prevents a path from revisiting an information set.
Given any behavioral continuation deviation, independently draw in
advance its action at each own information set. Every terminal suffix
then has the same product probability as under behavioral play.
The deviation's value is consequently a convex combination of values
of the pure continuation plans in \eqref{eq:sr}, and is no larger
than the prescribed value. Necessity is immediate because pure plans
are admissible. This proves full sequential rationality and the claim.
\end{proof}

All conditions are first-order ordered-field formulas. The radius
in \eqref{eq:consistency} is a real variable, not a quantifier over
sequences. Pure continuation plans are finite and can be enumerated;
alternatively behavioral deviations can be universally quantified using
the same polynomial suffix sums. Real quantifier elimination therefore
makes the exact relation effectively semialgebraic. This observation
alone is not an unrestricted efficiency claim.

\begin{lemma}[Nonempty equilibrium fibers]\label{lem:existence}
Every fixed admitted game has a sequential equilibrium.
\end{lemma}
\begin{proof}
When there are no strategic information sets, the empty assessment
suffices. Otherwise choose \(0<\varepsilon<1/\max_I|A(I)|\) and require
each action probability to be at least \(\varepsilon\).
Give each information set one decision agent with the original
player's unconditional utility. Parameterize its behavior by
\[
 \sigma_{I,a}=\varepsilon+
                  (1-|A(I)|\varepsilon)x_{I,a},
 \qquad x_I\text{ in the ordinary simplex}.
\]
An agent's payoff \(H_A(x)\) is continuous and affine in its own block.
To prove an agent equilibrium exists, put
\[
 g_{A,a}=\max\{0,H_A(a,x_{-A})-H_A(x)\},\qquad
 x_{A,a}\longmapsto
       \frac{x_{A,a}+g_{A,a}}{1+\sum_b g_{A,b}}.
\]
This continuous map of the compact convex product of simplexes has
a fixed point by Brouwer's theorem. If
\(G_A=\sum_b g_{A,b}>0\), the fixed-point equation gives
\(g_{A,a}=x_{A,a}G_A\). Every supported pure action would strictly
improve \(H_A(x)\), contradicting its expression as their weighted
average. Thus \(G_A=0\) for every agent.

Let \(\sigma^\varepsilon\) be the resulting constrained agent equilibrium
and \(\mu^\varepsilon\) its Bayes beliefs. Every information set has
positive reach. A unilateral agent change alters unconditional utility
by that reach probability times its conditional local payoff change.
The prescribed choice is therefore conditionally optimal among
\(\varepsilon\)-constrained distributions.

Perfect recall gives all nodes of an own information set \(J\) the same
remembered own sequence of earlier information sets and actions.
The product of own earlier action probabilities is common to those
nodes and cancels from Bayes' rule. With opponents and Nature fixed,
the conditional belief at \(J\) is independent of one's earlier
positive behavior. Moreover, if \(J\) follows \(I\) on one path,
all its nodes remember \(I\) and lie in its continuation.

Start with any constrained full continuation from \(I\). Replace its
own decisions by \(\sigma^\varepsilon\) in decreasing order of remembered
own-decision length. At a replacement at \(J\), later own decisions
already agree with \(\sigma^\varepsilon\), and the conditional belief
is \(\mu^\varepsilon_J\) by the cancellation just proved.
Local optimality makes the payoff change nonnegative; its contribution
at \(I\) is multiplied by the probability of reaching \(J\).
This proves optimality against every constrained continuation.

Let \(\varepsilon_m\downarrow0\) and take a convergent subsequence of
the compact assessment product. Its limit is consistent by construction.
For an arbitrary unconstrained continuation \(\delta_i\), use at each
own information set the constrained distribution
\(\varepsilon_m+(1-|A(J)|\varepsilon_m)\delta_{i,J,a}\).
It converges to \(\delta_i\). Passing to the limit in the finite
polynomial continuation inequalities proves \eqref{eq:sr}.
\end{proof}

The specially chosen perturbed equilibria in this existence proof
are not imposed on the witnesses of \eqref{eq:consistency}.

\begin{lemma}[Compact joint graph]\label{lem:compact}
The graph
\[
 \mathcal S=\{(\gamma,e):\gamma\in\Gamma,\
                   e=(\sigma,\mu),\ \SE(\gamma,e)\}
\]
is compact and semialgebraic, including at strategically off-path
information sets.
\end{lemma}
\begin{proof}
Semialgebraicity follows from Lemma~\ref{lem:exactse} and quantifier
elimination. Consider a convergent sequence
\((\gamma_m,\sigma_m,\mu_m)\in\mathcal S\) with limit
\((\gamma,\sigma,\mu)\). The weak polynomial inequalities
\eqref{eq:sr} persist at the limit.
For each \(m\), consistency supplies a completely mixed \(b_m\)
within \(1/m\) of \(\sigma_m\), whose Bayes beliefs at \(\gamma_m\)
are within \(1/m\) of \(\mu_m\).

Factor each prefix reach as
\[
 R_h(\gamma,b)=c_h(\gamma)a_h(b),
\]
where \(c_h\) is the product of Nature factors and \(a_h\) the product
of strategic factors. The finitely many \(c_h\) are positive and
continuous on \(\Gamma\). Consequently
\[
 \rho_m=\max_h\left|
     \frac{c_h(\gamma)}{c_h(\gamma_m)}-1\right|\longrightarrow0 .
\]
For \(\rho_m<1\), each reach weight at \(\gamma\) is its weight at
\(\gamma_m\) multiplied by a factor in
\([1-\rho_m,1+\rho_m]\). If the old normalized weights are \(v_h\)
and those factors are \(t_h\), their new normalized weights are
\(v_ht_h/\sum_jv_jt_j\). Their \(\ell^1\) difference is at most
\[
 \sum_h v_h
   \frac{|t_h-\sum_jv_jt_j|}{\sum_jv_jt_j}
       \le \frac{2\rho_m}{1-\rho_m}.
\]
This bound is independent of how small strategic reach becomes.
Hence the Bayes beliefs of \(b_m\) at the fixed \(\gamma\) converge
to \(\mu\), while \(b_m\to\sigma\). Consistency holds at the limit.
The graph is closed in the compact parameter/assessment product.
\end{proof}

The proof requires fixed uniformly positive Nature support.
It does not close a Bayes graph across a vanishing Nature-support
boundary. No such extension is used below.

\section{History recovery and the sharp informative fiber}\label{sec:fiber}
Assume \(M\) has full column rank and fix a known left inverse, for example
\begin{equation}\label{eq:leftinverse}
 A=(M^{\mathsf T}M)^{-1}M^{\mathsf T},\qquad AM=I_L.
\end{equation}
Then \(q=Ap\). A candidate population vector must satisfy \(Mq=p\)
and \(q\) must be a probability vector; otherwise it is incompatible.
For the present informative branch suppose \(q_z>0\) for every \(z\).

Let \(Z(h)\) be the terminal descendants of node \(h\), and put
\[
 r_h=\sum_{z\in Z(h)}q_z.
\]
At a Nature edge \(h\to ha\), recover
\begin{equation}\label{eq:naturerecovery}
 c^q_{h,a}=r_{ha}/r_h.
\end{equation}
Perfect recall makes each strategic information set \(I\) an antichain.
Let \(Z(I)\) be the disjoint union of \(Z(h)\), \(h\in I\), and
\(Z(I,a)\) the union of descendants of their action-\(a\) children.
Write \(r_I=\sum_{h\in I}r_h\) and \(r_{I,a}=\sum_{h\in I}r_{ha}\).
The candidate behavior and beliefs are
\begin{equation}\label{eq:behaviorrecovery}
 \sigma^q_{I,a}=\frac{r_{I,a}}{r_I},\qquad
 \mu^q_{I,h}=\frac{r_h}{r_I}.
\end{equation}
They are positive probability vectors, but behavioral compatibility
also requires
\begin{equation}\label{eq:compatibility}
 \frac{r_{ha}}{r_h}=\frac{r_{h'a}}{r_{h'}}
 \quad\text{for every }h,h'\in I,\ a\in A(I).
\end{equation}
These are testable polynomial equalities after cross multiplication.

\begin{lemma}[Recovery and its converse]\label{lem:recovery}
A positive \(q\) is generated by the stated information structure
and an admitted Nature vector exactly when \eqref{eq:compatibility}
holds and \(c^q\) is admitted. Nature probabilities, behavioral
probabilities, and consistent beliefs are then uniquely recovered
by \eqref{eq:naturerecovery}--\eqref{eq:behaviorrecovery}.
\end{lemma}
\begin{proof}
In any behavioral model, action probabilities agree at all nodes of
one information set, and conditional Nature probabilities equal the
primitive row. This proves necessity and uniqueness of the ratios.
Conversely start with root probability one and use each recovered
edge probability. At strategic nodes, \eqref{eq:compatibility}
makes that probability the common \(\sigma^q_{I,a}\).
Induction along the tree assigns node \(h\) probability \(r_h\),
because the factor on each edge is the child/parent mass ratio.
Thus the generated terminal law is exactly \(q\).
Bayes' rule gives \(\mu^q\). Every history and action has positive
probability, so consistent beliefs are unique; the constant completely
mixed strategy sequence supplies a consistency witness.
\end{proof}

The lemma has not assumed equilibrium. In particular, measurement
inversion alone does not establish information-set compatibility or
the shared Nature restrictions.

For an information set \(I\) owned by \(i\), define
\begin{align}
 m^0_{i,I,a}(q)&=
     \frac{\sum_{z\in Z(I,a)}q_zu^0_{i,z}}{r_{I,a}},&
 m^\varphi_{i,I,a}(q)&=
     \frac{\sum_{z\in Z(I,a)}q_z\varphi_{i,z}}{r_{I,a}}.
 \label{eq:conditionalfeatures}
\end{align}
Fix a reference action \(a_0(I)\). For each other action form a row
\begin{equation}\label{eq:Dandb}
 D_{i,I,a}=m^\varphi_{i,I,a}-m^\varphi_{i,I,a_0},\qquad
 b_{i,I,a}=m^0_{i,I,a_0}-m^0_{i,I,a}.
\end{equation}
Forced-action information sets add no row.

\begin{lemma}[Conditional features and full optimality]\label{lem:indifference}
At the recovered compatible profile, the value of choosing \(a\) at
\(I\) and then following the recovered continuation is
\begin{equation}\label{eq:Qconditional}
 Q_{i,I,a}(\beta)
    =m^0_{i,I,a}(q)+m^\varphi_{i,I,a}(q)\cdot\beta
    =\E_q[u_{i,Z}(\beta)\mid Z(I,a)].
\end{equation}
Its Bayes assessment is sequentially rational if and only if
\(D(q)\beta=b(q)\).
\end{lemma}
\begin{proof}
At every hidden node in \(I\), the probability of \(a\) is the same
positive \(\sigma^q_{I,a}\). Conditioning on that action leaves the
node distribution equal to \(\mu^q_I\). Later play is unchanged, so
the terminal conditional distribution is \(q\) restricted to \(Z(I,a)\).
This proves \eqref{eq:Qconditional}, with terminal utilities including
any already incurred costs under the fixed utility convention.

Sequential rationality implies
\(Q_{i,I,a}\le\sum_b\sigma^q_{I,b}Q_{i,I,b}\) for every \(a\).
All weights are positive, so all action values must coincide.
These are precisely the equations \(D(q)\beta=b(q)\).

Conversely suppose all those equations hold. Fix an arbitrary full
continuation deviation from \(I\), allowing zero action probabilities.
At every later own information set \(J\), perfect recall makes the
product of earlier own probabilities common to all its nodes.
When \(J\) is reached it cancels, leaving the original Bayes belief;
when it is not reached, it has no payoff contribution. Every node
of a later \(J\) lies in \(I\)'s continuation by perfect recall.
Replace deviating decisions by the prescribed behavior in decreasing
order of remembered own-decision length. At each replacement, later
behavior is already prescribed and all local action values coincide
by \eqref{eq:Qconditional}. The replacement changes no payoff.
After finitely many replacements the prescribed continuation remains.
Thus no original deviation improves the payoff. Consistency was
already established by Lemma~\ref{lem:recovery}.
\end{proof}

\begin{theorem}[Sharp informative identification]\label{thm:fiber}
If the inverse \(q=Ap\) is a positive probability vector, failure
of \(Mq=p\) or \eqref{eq:compatibility} makes the single-assessment
model incompatible. Otherwise its entire parameter identified set is
\begin{equation}\label{eq:sharpfiber}
 \Gamma_I(p)=\mathcal F(q)
   =\{(\alpha,\beta)\in\Gamma:
        B\alpha=c^q-c^0,\ D(q)\beta=b(q)\}.
\end{equation}
Every retained parameter has the recovered \((\sigma^q,\mu^q)\)
as an actual baseline sequential equilibrium with observation law \(p\).
\end{theorem}
\begin{proof}
Full column rank fixes \(q\) for every compatible assessment.
Lemma~\ref{lem:recovery} fixes its Nature vector, behavior, and beliefs;
Lemma~\ref{lem:indifference} imposes exactly the payoff equations.
This proves that every compatible parameter belongs to
\eqref{eq:sharpfiber}. Conversely each point in that polytope
supplies the admitted Nature vector and the compatible positive
profile, with consistent beliefs and full sequential optimality
by the same lemmas. Its terminal and observed laws are \(q,p\).
Both inclusions use actual assessments, proving sharpness.
\end{proof}

An empty polytope in \eqref{eq:sharpfiber} reports incompatibility
with the primitive restrictions. A zero terminal component does
not by itself report incompatibility: it places the problem in the
general exact branch of Section~\ref{sec:policysets}, retaining
strategically off-path histories and beliefs.

Let
\begin{equation}\label{eq:operator}
 \mathsf K(q)=\diag(B,D(q)).
\end{equation}
Full column rank is sufficient for point identification when the
fiber is nonempty. In general the fiber is a singleton exactly when
all its coordinate minima and maxima coincide; linear programs decide
this condition. If a true point lies in the relative interior of
\(\Gamma\), whose affine hull is \(E\gamma=f\), the singleton condition is
\begin{equation}\label{eq:relativecriterion}
 \ker\mathsf K(q)\cap\ker E=\{0\}.
\end{equation}
A nonzero vector in this intersection can be moved a small distance
in both signs within \(\Gamma\), preserving the data equations.
Conversely, a trivial intersection allows at most one point in
their affine-hull intersection. At a boundary, inequalities can make
the fiber a singleton without ambient full rank; the coordinate test
handles that case.

Observable primitive probabilities need not identify an arbitrary
structural index. Null directions of \(B\), additive utility shifts,
free scale changes expressible in the family, and payoff changes
invisible to continuation differences remain in \eqref{eq:sharpfiber}
unless joint primitive restrictions remove them.

Recovery uses finite sums, positive divisions, and linear algebra.
The feature formulas require terminal-event lists, not exponential
enumeration of pure continuation plans. With rational \(p\) and
rational primitives, the arrays have polynomial bit size in the
explicit tree, channel, features, and constraints. Standard rational
linear programming decides feasibility and coordinate extrema and
returns exact rational witnesses in polynomial bit complexity.
For algebraic data its representation must be specified; symbolic
population statements keep \(p\) as a free variable.

\section{Calibration, selection, and a nonempty informative family}
\label{sec:examples}
\begin{example}[An unknown channel is not calibration]
For two terminal labels let
\[
 M_\xi=
 \begin{pmatrix}(1+\xi)/2&(1-\xi)/2\\
                 (1-\xi)/2&(1+\xi)/2\end{pmatrix},
 \qquad 0<\xi\le1.
\]
Each matrix is invertible, but
\(p_1=1/2+\xi(q_1-1/2)\). Unknown \(\xi\) permits many pairs
\((\xi,q_1)\) with the same observed law. Known full rank without
known calibration and aligned labels is insufficient.
\end{example}

\begin{example}[Hidden equilibrium mixtures need not be one assessment]
In the simultaneous two-player coordination game, both players receive
one if their actions \(L,R\) match and zero otherwise. Its pure
equilibria \(LL,RR\) and its independent fair mixed equilibrium are
sequential equilibria in the standard perfect-recall tree.
Mix their selection probabilities as
\((1-\zeta)/2,(1-\zeta)/2,\zeta\), with \(0<\zeta<1\).
The resulting terminal law is positive and satisfies
\[
 q_{LL}=q_{RR}=(2-\zeta)/4,\qquad
 q_{LR}=q_{RL}=\zeta/4 .
\]
The second player's information set joins the nodes after \(L,R\).
Its conditional probability of choosing \(L\) is
\((2-\zeta)/2\) at the first node and \(\zeta/2\) at the second.
They differ, violating \eqref{eq:compatibility}. No single behavioral
assessment in that information structure generates this law.
Observed selector strata can instead be conditioned on separately.
\end{example}

\begin{proposition}[A primitive-guaranteed correlated Bayesian family]
\label{prop:bayesian}
For every positive integer \(m\), there is a nonempty compact class
with unknown correlated private-type law and unknown normalized payoff
coefficients, uniquely determined equilibrium behavior, and primitive
uniform lower bounds on all terminal probabilities and payoff rank.
\end{proposition}
\begin{proof}
There are two zero-sum players with private types \(t,s\in\{1,\ldots,m\}\).
Their unknown joint law is a matrix \(\mathsf Q\) with entries summing
to one and restrictions
\begin{equation}\label{eq:typeclass}
 \mathsf Q_{ts}\ge q_{\min}>0,\qquad
 \mathsf Q_{tt}\ge3\sum_{s\ne t}\mathsf Q_{ts}.
\end{equation}
This compact polytope is nonempty when
\(q_{\min}\le1/(4m^2)\), as witnessed by
\[
 \mathsf Q=\frac1m\left(\frac34 I_m+
                           \frac1{4m}\mathbf1\mathbf1^{\mathsf T}\right).
\]
It is invertible: if \(\mathsf Qx=0\), take \(t\) with
\(|x_t|=\|x\|_\infty\). Then
\(\mathsf Q_{tt}|x_t|
 \le\sum_{s\ne t}\mathsf Q_{ts}|x_s|
 \le\mathsf Q_{tt}|x_t|/3\), forcing \(x=0\).

After learning only their own types, the players simultaneously choose
zero or one. Player 1 receives \(\theta_t\) when both choose one,
one when both choose zero, and zero otherwise; player 2 receives
the negative. The unknown \(\theta_t\in[3/4,5/4]\) are calibrated
by the known zero and one payoffs.
Let \(p^a_t,q^a_s\) be their probabilities of action one, and put
\[
 K^{\rm type}_{ts}=\frac{\mathsf Q_{ts}}{\sum_s\mathsf Q_{ts}},
 \qquad r_t=\frac1{1+\theta_t}.
\]
Player 1's typewise indifference and player 2's typewise indifference
are respectively
\begin{equation}\label{eq:bayesindifference}
 K^{\rm type}q^a=r,\qquad
 \mathsf Q^{\mathsf T}
       \big((1+\theta)\mathbin{\odot}p^a-\mathbf1\big)=0 .
\end{equation}
Here \(\odot\) means coordinatewise multiplication. Invertibility gives
the only possible interior solution
\begin{equation}\label{eq:bayesbehavior}
 p^a=r,\qquad q^a=(K^{\rm type})^{-1}r .
\end{equation}

It is interior. The diagonal restrictions give
\(K^{\rm type}_{tt}\ge3/4\), so
\(\|K^{\rm type}-I\|_\infty\le1/2\).
The geometric matrix series gives
\(\|(K^{\rm type})^{-1}\|_\infty\le2\).
Since \(K^{\rm type}\mathbf1=\mathbf1\) and
\[
 r_t\in[4/9,4/7],\qquad
 \|r-\tfrac12\mathbf1\|_\infty\le1/14,
\]
we have \(\|q^a-\tfrac12\mathbf1\|_\infty\le1/7\).
Both players' action probabilities and their complements exceed \(1/3\).
Every conditional indifference holds. Each type has one genuine
decision, so the positive profile with its Bayes beliefs is a Bayesian
Nash and sequential equilibrium.

It is the unique equilibrium behavior. The ex ante game is bilinear
and zero-sum on the cubes of typewise probabilities. At the interior
opponent in \eqref{eq:bayesbehavior}, all own actions are indifferent.
Any other optimal \(p'\) must have zero payoff gradient in \(q^a\):
otherwise a small feasible perturbation of the interior opponent
would lower its guaranteed payoff below the saddle value.
The gradient equation in \eqref{eq:bayesindifference} and invertibility
of \(\mathsf Q\) imply \(p'=r\).
The symmetric argument with interior \(p^a\) gives
\(K^{\rm type}q'=r\), hence \(q'=q^a\).
Every Bayesian equilibrium is root saddle-optimal by the zero-sum
minimax inequalities; this argument does not posit general-sum uniqueness.

Every complete terminal \((t,s,a,b)\) consequently has probability
at least \(q_{\min}/9\). The correlated law \(\mathsf Q\) is recovered
by summing complete-history masses. In the payoff fiber, player 1's
type-\(t\) equation is
\[
 (K^{\rm type}q^a)_t\theta_t
       =1-(K^{\rm type}q^a)_t.
\]
Its coefficient matrix in \(\theta\) is diagonal with entries
\(r_t\ge4/9\), proving the claimed rank and conditioning.
\end{proof}

This family may have noisy post-play measurement. On its \(L\)
terminal labels take, for known rational \(0<\xi\le1\),
\begin{equation}\label{eq:noisychannel}
 M_\xi=\xi I_L+(1-\xi)L^{-1}\mathbf1\mathbf1^{\mathsf T}.
\end{equation}
It is stochastic, with eigenvalue one along \(\mathbf1\) and
\(\xi\) on its orthogonal complement. Its Euclidean operator inverse
norm is \(\|M_\xi^{-1}\|_2=1/\xi\).
It reveals no private type during play. Private types can also be
paths revealed in pieces before the terminal actions and persisting
through the earlier history. The general theorem itself concerns
arbitrary dynamic trees; this family demonstrates nonvacuous primitive
conditions, not a replacement of that theorem by one example.

\section{A statistical projection and parameter estimator}\label{sec:stats}
For the quantitative informative branch assume \(q_z\ge\lambda>0\).
Let \(m_c\) be the number of Nature entries,
\(J=\sum_I(|A(I)|-1)\) the number of contrast rows, and
\(R=m_c+J\). Let \(H_i\) be the maximum number of player \(i\)'s
decision nodes on a terminal path and \(H_*=\max_iH_i\), or zero
if there are no strategic decisions. Put \(a_M=\|A\|_2\).

Let \(\widehat p\) be the empirical cell-frequency vector and define
\begin{equation}\label{eq:projection}
 \Delta_\lambda=\{v\in\R^L:v_z\ge\lambda,\ \sum_zv_z=1\},
 \quad
 \widehat q=\underset{v\in\Delta_\lambda}{\operatorname{argmin}}\,
                 \|v-A\widehat p\|_2^2,\quad
 \Delta=\|\widehat q-q\|_1.
\end{equation}
The nonempty compact convex domain and strict convexity give a unique
projection. If \(L\lambda=1\), it is the constant vector \(\lambda\).
Otherwise its water-filling formula is
\[
 \widehat q_z=\max\{\lambda,(A\widehat p)_z-\tau\},
 \qquad \sum_z\widehat q_z=1 .
\]
Sort \(A\widehat p\), enumerate the active-coordinate counts, and solve
the last linear equation on each corresponding interval.
The feasible interval supplies \(\tau\); ties give the same vector.
Subtracting \(\tau\) at free coordinates and truncating the others
verifies the projection optimality inequalities. With rational inputs
this is an exact rational algorithm.

The statistic \(\widehat q\) need not satisfy
\eqref{eq:compatibility} or shared Nature restrictions.
It is never called an equilibrium law without those checks.

\begin{lemma}[Inversion risk and global conditioning]\label{lem:statbasic}
The projection satisfies
\begin{equation}\label{eq:S1}
 \E\Delta^2\le\frac{La_M^2}{n},\qquad
 \E\Delta\le\frac{\sqrt L\,a_M}{\sqrt n},\qquad \Delta\le2 .
\end{equation}
If probability vectors \(v,w\) and an event \(E\) satisfy
\(v(E)\ge\lambda\), \(w(E)>0\), and \(|f|\le F_0\), then
\begin{equation}\label{eq:S2}
 |\E_v[f\mid E]-\E_w[f\mid E]|
          \le\frac{2F_0}{\lambda}\|v-w\|_1.
\end{equation}
For nonnegative vectors \(x,y\) with positive sums \(s,t\),
\begin{equation}\label{eq:S3}
 \|x/s-y/t\|_1\le 2\|x-y\|_1/s .
\end{equation}
\end{lemma}
\begin{proof}
Euclidean projection is nonexpansive. Indeed, if \(x=\Pi(v)\) and
\(y=\Pi(w)\), optimality gives
\(\langle v-x,y-x\rangle\le0\) and
\(\langle w-y,x-y\rangle\le0\). Adding yields
\[
 \|x-y\|_2^2\le\langle x-y,v-w\rangle
                    \le\|x-y\|_2\|v-w\|_2 .
\]
Because \(q=Ap\in\Delta_\lambda\),
\[
 \|\widehat q-q\|_2\le a_M\|\widehat p-p\|_2,\qquad
 \Delta\le\sqrt L\,a_M\|\widehat p-p\|_2 .
\]
The multinomial coordinate variances give
\(\E\|\widehat p-p\|_2^2=(1-\|p\|_2^2)/n\le1/n\).
Cauchy--Schwarz gives the first two assertions of \eqref{eq:S1};
the last follows because both vectors are probability vectors.

For \eqref{eq:S2}, put \(b=v(E)\), \(b'=w(E)\),
\(t_0=\sum_Ev_zf_z\), and \(t'_0=\sum_Ew_zf_z\).
The exact identity
\[
 \frac{t_0}{b}-\frac{t'_0}{b'}
   =\frac{t_0-t'_0}{b}
      +\frac{t'_0}{b'}\,\frac{b'-b}{b}
\]
and \(|t'_0/b'|\le F_0\) prove the bound.
Only the reference denominator requires a lower bound.
Finally
\[
 x/s-y/t=(x-y)/s+(y/t)(t-s)/s,\qquad
 |t-s|\le\|x-y\|_1
\]
proves \eqref{eq:S3}. Both conditional bounds are global; no
small-error event is needed.
\end{proof}

Every nonempty prefix or information-action event contains a terminal,
so its \(q\)- and \(\widehat q\)-mass is at least \(\lambda\).
Define \(c^{\rm raw}_{h,a}=\widehat q(Z(ha))/\widehat q(Z(h))\)
at Nature nodes. Define \(\widehat D,\widehat b\) by
\eqref{eq:conditionalfeatures}--\eqref{eq:Dandb} with \(\widehat q\)
in place of \(q\). They are rational conditional-moment statistics,
whether or not \(\widehat q\) is behaviorally compatible.
Lemma~\ref{lem:statbasic} gives
\begin{equation}\label{eq:S4}
 \begin{split}
 \|c^{\rm raw}-c^q\|_\infty&\le2\Delta/\lambda,\\
 \max_{j,k}|(\widehat D-D)_{j,k}|&\le4F\Delta/\lambda,\\
 \|\widehat b-b\|_\infty&\le4U_0\Delta/\lambda .
 \end{split}
\end{equation}
Uniformly for every \(\beta\) admitted by \(\Gamma\),
\begin{equation}\label{eq:S5}
 \|(\widehat D-D)\beta-(\widehat b-b)\|_\infty
                    \le4U\Delta/\lambda .
\end{equation}
This follows from the entry bounds and \(\|\beta\|_1\le B_1\),
or by applying \eqref{eq:S2} directly to \(u_{i,Z}(\beta)\).

Put \(C=\max\{2,4U\}/\lambda\). The estimator solves the rational LP
\begin{equation}\label{eq:fitLP}
 \begin{aligned}
 \text{minimize }&t\\
 \text{over }&(\alpha,\beta)\in\Gamma,\quad t\ge0\\
 \text{subject to }&
   \|B\alpha-(c^{\rm raw}-c^0)\|_\infty\le t,\qquad
   \|\widehat D\beta-\widehat b\|_\infty\le t .
 \end{aligned}
\end{equation}
Its optimizer \((\widehat\gamma,\widehat t)\) exists because the
minimized residual is continuous on compact \(\Gamma\).
A deterministic exact rational choice is obtained by optimizing \(t\)
and then successively minimizing coordinates on its optimizer face.

\begin{theorem}[Every-sample fit and conditioned risk]\label{thm:risk}
For every point \(\gamma\) of the entire true fiber \(\mathcal F(q)\),
\begin{align}
 \widehat t&\le C\Delta, \label{eq:S6}\\
 \|\mathsf K(q)(\widehat\gamma-\gamma)\|_\infty
       &\le2C\Delta. \label{eq:S7}
\end{align}
If \(\mathsf K(q)\) has full column rank with smallest singular
value \(\kappa>0\), the fiber is a singleton \(\gamma_0\), and
for any known \(G\ge\diam_2(\Gamma)\),
\begin{equation}\label{eq:S8}
 \E\|\widehat\gamma-\gamma_0\|_2^2
 \le\min\left\{G^2,\frac{4RC^2La_M^2}{\kappa^2n}\right\}.
\end{equation}
\end{theorem}
\begin{proof}
Every fiber point has residual at most \(C\Delta\) in
\eqref{eq:fitLP}, proving \eqref{eq:S6}.
The fitted chance block differs from the true chance block by at
most \(\widehat t+2\Delta/\lambda\).
The fitted payoff block satisfies
\(\|D\widehat\beta-b\|_\infty
 \le\widehat t+4U\Delta/\lambda\), since
\(\widehat\gamma\in\Gamma\) permits \eqref{eq:S5}.
Both are at most \(2C\Delta\). Subtracting any fiber point's exact
equations proves \eqref{eq:S7}.
Under full rank,
\(\|\widehat\gamma-\gamma_0\|_2\le2\sqrt R\,C\Delta/\kappa\).
Square, use \eqref{eq:S1}, and separately use both points'
membership in \(\Gamma\) to obtain \eqref{eq:S8}.
\end{proof}

The bound is unconditional, without an estimated inverse,
a logarithm of \(n\), or a fixed-probability remainder.
A valid \(G\) is the square root of the sum of squared coordinate
range widths of \(\Gamma\), found by \(2d\) LPs.
Computing its exact Euclidean diameter is not claimed to be an LP.
Under rank deficiency, \eqref{eq:S7} estimates the observable linear
combinations; it does not identify an arbitrary chosen point of the fiber.

\section{Whole-fiber confidence and diameter certificates}
\label{sec:confidence}
Fix a confidence error \(0<\eta<1\), distinct from the chance parameter
\(\alpha\). The elementary coordinate Hoeffding bound gives, with
probability at least \(1-\eta\),
\[
 \|\widehat p-p\|_\infty\le h_n,\qquad
 h_n=\sqrt{\frac{\log(2K/\eta)}{2n}}.
\]
For completeness, the logarithm of the centered moment-generating
function of a Bernoulli variable has value and derivative zero at
zero, and second derivative equal to a tilted Bernoulli variance,
at most \(1/4\). Twice integrating yields
\(\log\E e^{t(X-\E X)}\le t^2/8\) for every real \(t\).
Independence and exponential Markov, optimized at \(t=4h\),
give each upper tail \(e^{-2nh^2}\); negative \(t\) gives the
same lower-tail bound. A union bound over both tails and all \(K\)
coordinates proves the assertion.

Define
\begin{equation}\label{eq:S9}
 e_n=\min\{2,\sqrt{LK}\,a_Mh_n\}.
\end{equation}
The same event gives \(\Delta\le e_n\).
Our outer confidence polytope is
\begin{equation}\label{eq:confidencepoly}
 \begin{split}
 \mathcal C_n=\{(\alpha,\beta)\in\Gamma:\;&
  \|B\alpha-(c^{\rm raw}-c^0)\|_\infty\le2e_n/\lambda,\\
 &\|\widehat D\beta-\widehat b\|_\infty\le4Ue_n/\lambda\}.
 \end{split}
\end{equation}

\begin{theorem}[Finite-sample coverage of the entire fiber]
\label{thm:confidence}
Uniformly over the declared quantitative law class,
\begin{equation}\label{eq:S10}
 \Prob\{\mathcal F(q)\subseteq\mathcal C_n\}\ge1-\eta .
\end{equation}
This remains valid at rank-deficient chance/payoff operators.
\end{theorem}
\begin{proof}
On \(\Delta\le e_n\), \eqref{eq:S4}--\eqref{eq:S5} imply both
inequalities in \eqref{eq:confidencepoly} simultaneously for every
point of the exact fiber. The implication is deterministic on the
single concentration event; no parameter or equilibrium is selected.
\end{proof}

A fully rational implementation can replace \(e_n\) everywhere by
a conservative rational radius. For rational \(\eta\), let \(b_n\)
be the least nonnegative integer with \(2^{b_n}\ge2K/\eta\), and put
\begin{equation}\label{eq:rationalradius}
 \bar h_n=\frac{\lceil\sqrt{2nb_n}\rceil}{2n},\qquad
 g=\lceil\sqrt{LK}\rceil,\qquad
 \bar e_n=\min\{2,g\bar a_M\bar h_n\},
\end{equation}
where rational \(\bar a_M\ge a_M\) can upper-bound the Frobenius norm
of \(A\). Since \(\log(2K/\eta)\le b_n\) and
\(\bar h_n\ge\sqrt{b_n/(2n)}\), we have \(\bar e_n\ge e_n\).
Integer comparisons and square roots construct it. At fixed input
and confidence it is \(O(n^{-1/2})\), with an additional
\(O(n^{-1})\) rounding term. All subsequent formulas permit this
replacement. The radius cap uses the deterministic \(\Delta\le2\).

The polytope can be empty outside the coverage event and is reported
as empty. It is not silently replaced by a point estimate.
An optional replacement by \(\Gamma\) is conservative but is not
used in the expected-diameter proof. At radius two,
\(\mathcal C_n=\Gamma\): the compared chance entries are probabilities
and every payoff contrast has magnitude at most \(2U\).

For a nonempty \(\mathcal C_n\), obtain each coordinate's minimum
\(l_j\) and maximum \(u_j\) with \(2d\) LPs. Then
\[
 \diam_2(\mathcal C_n)\le
             \sqrt{\sum_{j=1}^d(u_j-l_j)^2}.
\]
With the rational radius these endpoints are exact rational extrema.
The coordinate ranges are sharp for this outer polytope; the resulting
bounding-box diameter need not equal its Euclidean diameter.

Let \(\widehat{\mathsf K}=\diag(B,\widehat D)\).
If its smallest singular value \(\widehat\kappa\) is positive,
the difference of any two candidate vectors obeys the rowwise
tolerances in \eqref{eq:confidencepoly}, giving
\begin{equation}\label{eq:S11}
 \begin{split}
 \diam_2(\mathcal C_n)
 &\le\min\left\{G,
  \frac{2e_n\sqrt{m_c(2/\lambda)^2+J(4U/\lambda)^2}}
       {\widehat\kappa}\right\}\\
 &\le\min\{G,2\sqrt R\,Ce_n/\widehat\kappa\}.
 \end{split}
\end{equation}
A certified positive lower bound for \(\widehat\kappa\) suffices.
An estimated rank test is unnecessary for coverage and does not,
by itself, establish population point identification.

\begin{proposition}[Every-sample and expected diameter]
\label{prop:diameter}
Suppose the true \(\mathsf K(q)\) has smallest singular value
\(\kappa>0\). With \(\diam(\varnothing)=0\),
\begin{align}
 \diam_2(\mathcal C_n)
  &\le\min\{G,2\sqrt R\,C(e_n+\Delta)/\kappa\},
                                      \label{eq:S12}\\
 \E\diam_2(\mathcal C_n)
  &\le\min\{G,2\sqrt R\,C(e_n+\sqrt L\,a_M/\sqrt n)/\kappa\}.
                                      \label{eq:S13}
\end{align}
\end{proposition}
\begin{proof}
Fix any true fiber point \(\gamma_0\). For any
\(\gamma=(\alpha,\beta)\in\mathcal C_n\), the chance block differs
from the true one by at most \(2(e_n+\Delta)/\lambda\).
For its payoff block, \eqref{eq:S5} applies to this candidate
\(\beta\in\Gamma\), giving a difference at most
\(4U(e_n+\Delta)/\lambda\). Thus
\[
 \|\mathsf K(q)(\gamma-\gamma_0)\|_\infty
                    \le C(e_n+\Delta).
\]
Every candidate is within
\(\sqrt R\,C(e_n+\Delta)/\kappa\) of \(\gamma_0\), whether or not
\(\gamma_0\) belongs to this sample's confidence polytope.
Triangle inequality proves \eqref{eq:S12}; the empty set satisfies
it by convention. Take expectations, use \eqref{eq:S1}, and retain
the independent \(G\) cap to obtain \eqref{eq:S13}.
\end{proof}

There is no \(\eta G\) remainder. Fixed positive conditioning gives
an \(O(n^{-1/2})\) expected diameter at fixed confidence. No point-
shrinking claim is made when a nontrivial true fiber remains.
For reference, \eqref{eq:S4} also implies
\[
 \|\widehat{\mathsf K}-\mathsf K(q)\|_2
          \le (4F\sqrt{Jd_\beta}/\lambda)\Delta .
\]
This Frobenius bound can transfer a supplied population rank margin
on a concentration event; neither the estimator risk nor
\eqref{eq:S12} requires such a transfer.

\section{An actual assessment and every-continuation certificates}
\label{sec:actual}
Use one behavior vector per information set,
\begin{equation}\label{eq:actualconstruction}
 \widehat\sigma_{I,a}
   =\frac{\widehat q(Z(I,a))}{\widehat q(Z(I))},\qquad
 c^\star=c^0+B\widehat\alpha,\qquad
 u^\star_{i,z}=u_{i,z}(\widehat\beta).
\end{equation}
The action events partition \(Z(I)\), so this is a proper distribution.
Its numerator is at least \(\lambda\) and denominator at most one;
hence \(\widehat\sigma_{I,a}\ge\lambda\).
Let \(q^\star,r^\star_h\) be the actual terminal and prefix
probabilities obtained by multiplying \(c^\star,\widehat\sigma\)
along the original tree, and define
\begin{equation}\label{eq:actualbelief}
 \mu^\star_{I,h}
     =\frac{r^\star_h}{\sum_{h'\in I}r^\star_{h'}}.
\end{equation}

Because \(\widehat\gamma\in\Gamma\), all shared Nature restrictions
are exactly satisfied. Let \(c_{\min}>0\) be the minimum included
Nature entry over \(\Gamma\), computed by finitely many LPs;
take \(c_{\min}=1\) if there are no Nature edges.
Every terminal has probability at least
\(\min\{c_{\min},\lambda\}^{H}>0\).
Thus \(\mu^\star\) is defined and satisfies Bayes' rule exactly.
The constant completely mixed strategy sequence proves consistency.
In general \(q^\star\ne\widehat q\); node-normalized
\(\widehat q\) weights are not used as these beliefs.

Equations \eqref{eq:S4} and \eqref{eq:S6} give
\[
 \|c^\star-c^q\|_\infty\le\widehat t+2\Delta/\lambda
                         \le2C\Delta,\qquad
 \|\widehat\sigma-\sigma^q\|_\infty\le2\Delta/\lambda\le C\Delta.
\]
Every edge factor therefore differs by at most
\(\delta_{\rm edge}=2C\Delta\).
For factors in \([0,1]\), telescoping products gives
\[
 \left|\prod_{j=1}^s x_j-\prod_{j=1}^s y_j\right|
                    \le\sum_{j=1}^s|x_j-y_j|.
\]
Each path has at most \(H\) factors. Consequently
\begin{equation}\label{eq:S14}
 |q^\star_z-q_z|\le H\delta_{\rm edge},\qquad
 \|q^\star-q\|_1\le2LHC\Delta .
\end{equation}
The last right-hand side may always be capped by two.
The same prefix-product bound and \eqref{eq:S3}, using true
information-set reach at least \(\lambda\), yield
\begin{equation}\label{eq:S15}
 \|\mu^\star_I-\mu^q_I\|_1
 \le\frac{2|I|H\delta_{\rm edge}}{\lambda}
 \le\frac{4N_{\rm tr}HC\Delta}{\lambda}.
\end{equation}
This is a global normalization bound, requiring no small-error split
or estimated reach lower bound.

All following action values concern this \emph{actual} reconstructed
game. The common positive action probability at \(I\) cancels in
conditioning, so
\begin{equation}\label{eq:S16}
 Q^\star_{i,I,a}
    =\E_{q^\star}[u_{i,Z}(\widehat\beta)\mid Z(I,a)],\qquad
 V^\star_{i,I}=\sum_a\widehat\sigma_{I,a}Q^\star_{i,I,a}.
\end{equation}
Define the exact local certificate
\[
 \ell_{\rm loc}=\max_{i,I,a}[Q^\star_{i,I,a}-V^\star_{i,I}]_+,
\]
or zero if no player moves. It is computable as a rational number
under the rational implementation.

The true-\(q\) contrasts at \(\widehat\beta\) are the entries of
\(D(q)\widehat\beta-b(q)\), bounded by \(2C\Delta\).
Apply \eqref{eq:S2} with reference \(q\), other law \(q^\star\),
and payoff bounded by \(U\), separately for two actions. With
\eqref{eq:S14} this gives
\begin{equation}\label{eq:S17}
 |Q^\star_{i,I,a}-Q^\star_{i,I,a_0}|
       \le C_{\rm loc}\Delta,\qquad
 C_{\rm loc}=2C+\frac{8ULHC}{\lambda}.
\end{equation}
Two arbitrary actions differ by at most \(2C_{\rm loc}\Delta\);
the prescribed value is a convex combination, whence
\begin{equation}\label{eq:S18}
 \ell_{\rm loc}\le2C_{\rm loc}\Delta .
\end{equation}

\begin{lemma}[Quantitative backward replacement]\label{lem:backwardbound}
In any finite perfect-recall game with positive Nature probabilities,
completely mixed behavior, and its exact Bayes beliefs, let
\(\ell_{\rm loc}\) be the local certificate just defined using
that game. For every player \(i\), starting information set \(I\),
and full continuation deviation \(\delta_i\),
\begin{equation}\label{eq:S19}
 V_{i,I}(\delta_i)-V_{i,I}(\text{prescribed continuation})
                         \le H_i\ell_{\rm loc}.
\end{equation}
The deviation can randomize, use zero action probabilities, and depend
on the full information histories allowed by the tree.
\end{lemma}
\begin{proof}
At an own information set \(J\), perfect recall fixes the remembered
sequence of own earlier information sets and actions. The product
of one's earlier behavioral probabilities is common at all nodes
of \(J\). Replacing that earlier behavior multiplies their reach
weights by the same factor. A positive factor cancels in conditional
beliefs; a zero factor makes \(J\) unreachable.
Thus any reached \(J\) under a deviation has the original Bayes belief.
The same argument applies to play started conditionally at \(I\):
a later own \(J\) remembers \(I\) and the action there, so all its
nodes belong to that continuation, and the common factor cancels.

Order own information sets after \(I\), including \(I\), by decreasing
remembered own-decision length. Replace deviating choices by the
prescribed ones in this order. When replacing \(J\), all later own
choices are prescribed. If \(J\) is reached, its belief is the one
used to compute \(Q_{i,J,a}\). The current deviating mixture can
exceed the prescribed mixture value by at most \(\ell_{\rm loc}\).
The reduction in payoff from replacing it is therefore at most
\(\ell_{\rm loc}\) times the probability of reaching \(J\) from \(I\).
An unreached set contributes zero.

Earlier own decisions are still the original deviation's decisions.
Replacements later than \(J\) cannot affect its reach, and sets with
the same remembered depth cannot follow one another.
All these reach probabilities are therefore those under the
\emph{same original} deviation from \(I\).
Their sum is its expected number of own decision nodes visited,
at most \(H_i\). Summing the payoff changes telescopically proves
\eqref{eq:S19}. This proves the full continuation bound rather than
assuming a one-step-deviation principle.

Behavioral deviations are covered directly. For privately randomized
complete plans, condition on the private randomization, apply the
bound to each deterministic plan, and average.
\end{proof}

Let \(\Reg_{\SE}\) be the maximum positive full continuation improvement
over all players and information sets, evaluated with \(\mu^\star\)
in the actual reconstructed game. Then
\begin{equation}\label{eq:S20}
 \begin{split}
 \Reg_{\SE}(\widehat\gamma,\widehat\sigma,\mu^\star)
   &\le H_*\ell_{\rm loc}\le C_{\rm eq}\Delta,\\
 C_{\rm eq}&=2H_*C_{\rm loc}
       =4H_*C(1+4ULH/\lambda).
 \end{split}
\end{equation}
Equation \eqref{eq:S1} and the concentration event give
\begin{align}
 \E\Reg_{\SE}^2
       &\le\min\{4U^2,C_{\rm eq}^2La_M^2/n\},\label{eq:S21}\\
 \Prob\{\Reg_{\SE}\le\min\{2U,C_{\rm eq}e_n\}\}
       &\ge1-\eta .\label{eq:S22}
\end{align}
Both utility comparisons lie in \([-U,U]\), which justifies the caps.
The rational radius can replace \(e_n\).

More strongly, \(H_*\ell_{\rm loc}\) is an a posteriori certificate
valid on every data set without knowledge of \(q,\Delta\), or the
true parameter. Primitive feasibility, common information-set
behavior, and Bayes consistency are exact. Only sequential optimality
has a residual. If \(\ell_{\rm loc}=0\), the assessment is an exact
sequential equilibrium; otherwise that assertion is not made.

\section{Effective statistical computation and its limits}
\label{sec:statcomputation}
For rationally encoded \(M,c^0,B,u^0,\varphi,\lambda,B_1,\Gamma,\eta\),
the following procedure is finite and explicit:
\begin{enumerate}
\item Compute \(A\) by exact rational linear algebra, form the
cell frequencies, and compute \(\widehat q\) by sorting and water filling.
\item Sum prefix and information-action masses and features, then form
\(c^{\rm raw},\widehat D,\widehat b\).
\item Solve \eqref{eq:fitLP} over the full joint \(\Gamma\);
construct \(\bar e_n\) and \(\mathcal C_n\).
\item Decide confidence-polytope feasibility and, when nonempty,
compute its coordinate extrema and certified diameter bound.
\item Form \(c^\star,\widehat\sigma,q^\star,\mu^\star\), compute actual
action values, and report the certificate \(H_*\ell_{\rm loc}\).
\end{enumerate}

There are polynomially many rational operations and LPs in the
explicit tree, channel, utility features, and parameter inequalities.
Sorting uses \(O(L\log L)\) comparisons.
Even naively listing the terminal incidence of each event has
polynomial explicit-input size.
Empirical event denominators are at least \(\lambda\).
Actual Bayes denominators are positive, possibly as small as
\(\min\{c_{\min},\lambda\}^{H}\); finite path products and sums
still have polynomial rational bit length in the explicit tree
and input bit length. Small probability is not large representation
length by itself.

Standard exact rational linear algebra has polynomial bit complexity.
We use the classical polynomial-bit algorithm for rational linear
programming, with explicitly listed rational constraints of polynomial
bit length. Deterministic optimizer selection takes at most \(d\)
additional LPs. Integer comparisons and square roots construct
\eqref{eq:rationalradius}; a rational sum of squares supplies a
certified upper bound for \(\|A\|_{\rm F}^2\).
Spectral computation is not needed for coverage, the estimator,
or the coordinate-range certificate. To report a spectral diameter
bound, a rational positive-semidefiniteness certificate for
\(\widehat{\mathsf K}^{\mathsf T}\widehat{\mathsf K}
 -\kappa_{\rm low}^2I_d\) supplies a valid positive
\(\kappa_{\rm low}\).

The procedure is polynomial in the explicit input and in \(\log n\)
when cell counts are supplied. Reading \(n\) raw observations costs
at least \(n\) operations. This is neither a strongly polynomial LP
claim nor a polynomial claim in succinct horizon \(T\).
Algebraic data need their exact representation and corresponding
arithmetic guarantees; arbitrary unencoded real input is excluded
from finite-bit assertions.

The population fiber is sharp. Its confidence polytope is a conservative
outer set, not another sharp population fiber.
Finite-sample incompatibility of the raw projection with
\eqref{eq:compatibility} is expected, not an exact rejection test.
An empty confidence polytope, in contrast, has probability at most
\(\eta\) under the declared model because the nonempty true fiber
is contained on the coverage event.

The quantitative constants expose measurement conditioning \(a_M\),
history conditioning \(\lambda\), payoff bounds, and the identification
margin \(\kappa\). They do not remain bounded as those conditions
degenerate. No optimal dependence on these constants is claimed.
Unknown calibration or a rank-deficient channel does not inherit
the inversion rate. Without full history support the generic exact
method remains valid, but these ratio rates require another analysis.
A nontrivial payoff/chance fiber remains set-valued.

\section{Sharp general policy sets and comparative selection}
\label{sec:policysets}
Return to the generic class, permitting rank-deficient \(M\) and
strategically zero terminal probabilities. Let \(\mathcal S_0\)
be the baseline equilibrium graph and \(\mathcal S_\ell\) the
policy graph from Section~\ref{sec:se}. Let
\(O(\gamma,\sigma)=Mq(\gamma,\sigma)\). For population law \(p\), define
\begin{equation}\label{eq:policysets}
 \begin{aligned}
 \mathcal A(p)&=\{(\gamma,e^0)\in\mathcal S_0:
                         O(\gamma,\sigma^0)=p\},\\
 \Gamma_I(p)&=\operatorname{proj}_\gamma\mathcal A(p),\\
 \mathcal J_\ell(p)&=\{(\gamma,e^0,e^\ell):
        (\gamma,e^0)\in\mathcal A(p),\
        (\gamma,e^\ell)\in\mathcal S_\ell\},\\
 \mathcal W_I(\ell;p)&=
   \{W_\ell(\gamma,\sigma^\ell):
                       (\gamma,e^0,e^\ell)\in\mathcal J_\ell(p)\}.
 \end{aligned}
\end{equation}

\begin{theorem}[Exact ranges with attained scalar endpoints]
\label{thm:generalsets}
These are exactly the canonical parameter set and all-selection
policy range. They are compact semialgebraic sets, possibly empty.
If \(\Gamma_I(p)\ne\varnothing\), each policy range is nonempty and
its scalar minimum and maximum are attained.
The full range need not be an interval.
\end{theorem}
\begin{proof}
Lemma~\ref{lem:exactse} makes every feasible tuple an actual baseline
SE producing \(p\) and an actual policy SE at the same parameter.
Conversely every canonical compatible model supplies such a tuple.
This proves sharpness in both directions, not merely containment.
Lemma~\ref{lem:compact} and the closed observation equality make
\(\mathcal A(p)\) compact. Its projection and the shared-parameter
product defining \(\mathcal J_\ell(p)\) are compact.
Continuity of \eqref{eq:policyoutcome} preserves compactness of
its image; polynomial quantification preserves semialgebraicity.
Nonempty policy fibers follow from Lemma~\ref{lem:existence},
and continuous scalar functions attain extrema on compact sets.
None of these operations fills gaps in a disconnected image.
\end{proof}

Without a supplied cross-policy selector, a sharp same-parameter
adverse comparison range is
\begin{equation}\label{eq:pairedpolicy}
 \begin{split}
 \{W_\ell(\gamma,\sigma^\ell)-W_{\ell'}(\gamma,\sigma^{\ell'}):\
 &(\gamma,e^0)\in\mathcal A(p),\\
 &(\gamma,e^\ell)\in\mathcal S_\ell,\
   (\gamma,e^{\ell'})\in\mathcal S_{\ell'}\}.
 \end{split}
\end{equation}
It permits every pair of policy equilibria, retaining the same
structural parameter. Its endpoints are attained by the same argument.
Subtracting marginal endpoints can use different parameters and
is only an explicitly broader bound. A supplied closed semialgebraic
joint selection relation can instead be intersected with
\eqref{eq:pairedpolicy}; nonemptiness must then be checked.
No such relation is assumed by default.

\section{Exact computation, endpoint errors, and coverage transfer}
\label{sec:qe}
All defining constraints above are explicit finite first-order formulas.
Classical effective real-closed-field quantifier elimination
\cite{Basu} gives descriptions of the full \(\Gamma_I\) and
\(\mathcal W_I\), including disconnected parts.
Emptiness is checked first. Rational and represented algebraic
numerical inputs are admissible; alternatively \(p\) is left symbolic.

For a nonempty feasible tuple set \(\mathcal J\), its scalar extrema
are specified by
\begin{equation}\label{eq:endpointformula}
 \begin{aligned}
 &\exists y\in\mathcal J:\ W(y)=w_{\min},\qquad
   \forall y\in\mathcal J:\ W(y)\ge w_{\min},\\
 &\exists y\in\mathcal J:\ W(y)=w_{\max},\qquad
   \forall y\in\mathcal J:\ W(y)\le w_{\max}.
 \end{aligned}
\end{equation}
Quantifier elimination and real-root isolation compute these algebraic
numbers. Semialgebraic sample-point extraction supplies algebraic
tuples attaining them, including actual baseline and policy SE
assessments. Their incentive residuals and consistency constraints
are exact; they are not only numerical best-response tests.

A rational bound \(|W|\le B_W\) is computable.
First bound the coordinates of compact \(\Gamma\) by rational LPs;
then bound each terminal polynomial by the sum of its absolute
coefficient-monomial bounds. Expected \(W\) is a convex combination
of those terminal values.
Start from rational \([L_{\rm init},U_{\rm init}]\), for instance
\([-B_W,B_W]\), and any feasible tuple.
For the maximum, query whether there exists a feasible tuple with
\(W\ge q_{\rm mid}\). A true answer gives a lower bound and a
witness; a false answer gives an upper bound. For the minimum use
\(W\le q_{\rm mid}\), with reversed interval updates.
After
\begin{equation}\label{eq:bisectcount}
 \max\left\{0,\left\lceil
       \log_2\frac{U_{\rm init}-L_{\rm init}}{\varepsilon}
                     \right\rceil\right\}
\end{equation}
queries when the initial width is positive, the interval width is
at most \(\varepsilon\). Zero width requires no queries.
For the maximum, the saved witness has value at least the maintained
lower bound, even if no midpoint query was true because the initial
witness is at least \(L_{\rm init}\). It is within \(\varepsilon\)
of the global maximum. The minimum argument is the reverse.
The certificate records the formulas, initial bounds, true-query
witnesses, and exact real-algebraic decision certificates for false
queries. This is a complete algorithm, not a presumed global oracle.

Unrestricted quantifier elimination may be expensive.
When \emph{all} compiled variable counts, degrees, and structural
tree bounds are fixed, the operation bound and integer bit-size clause
of Basu's Theorem 2.27 \cite[p.~16]{Basu} give polynomial bit cost
in dense coefficient/formula size and rational precision-input length.
Theorem 2.4 \cite[p.~6]{Basu} gives the usual CAD arithmetic bound.
All belief, tremble, continuation, and primitive variables must be
counted. This claim is distinct from polynomial explicit-tree LP
computation and implies no polynomial succinct-horizon algorithm.

For generic iid sampling, with probability at least \(1-\eta\),
\(\|\widehat p-p\|_\infty\le h_n\), by the Bernoulli argument
of Section~\ref{sec:confidence}. Replace the baseline observation
equality by
\[
 \|O(\gamma,\sigma^0)-\widehat p\|_\infty\le h_n,
\]
or a rational conservative radius. On that event, every tuple in
\(\mathcal A(p)\) remains feasible with its original exact SE witness.
Projection covers the whole parameter set; attaching every exact
policy equilibrium covers every policy range, simultaneously on the
same event. The same reasoning covers \eqref{eq:pairedpolicy}.
This generic coverage argument is retained from \cite{Prior}.

More generally, let \(E_n\) be a measurable event of probability at
least \(1-\eta\) on which an outer parameter set \(\mathcal C_n\)
contains the whole identified set. Define
\[
 \mathcal V_n(\ell)=
 \{W_\ell(\gamma,\sigma^\ell):
          \gamma\in\mathcal C_n\cap\Gamma,\
          (\gamma,e^\ell)\in\mathcal S_\ell\}.
\]
Set inclusion gives the following common-event assertion, writing
\(\omega\) for the sampled data:
\begin{equation}\label{eq:S23}
 \Prob(E_n)\ge1-\eta,\qquad
 \forall\omega\in E_n\ \forall\ell:\quad
 \mathcal W_I(\ell;p)\subseteq\mathcal V_n(\ell;\omega).
\end{equation}
It holds for every policy on one measurable event, without needing
measurability of an uncountable intersection of policy-specific events.
Thus the informative confidence polytope can be used without solving
the baseline equilibrium problem again.
Use the identical joint selection relation on both sides if one is
part of the target. Exact policy ranges over the outer polytope can
be computed by the preceding algorithm.
Certified outer supersets preserve coverage; a small equilibrium
residual alone does not imply outer containment.
The approximate baseline assessment in Section~\ref{sec:actual}
is never substituted for an exact policy equilibrium here.

\section{Zero-sum policies: realization flows and a value LP}
\label{sec:sequences}
For the faster policy branch there are two strategic players and
Nature, perfect recall, and \(u_{2,z}=-u_{1,z}\).
The target \(W_\ell\) is precisely player 1's expected zero-sum
utility, not an arbitrary action, allocation, revenue, or welfare
statistic. The following classical sequence construction
\cite[Section~2, pp.~249--250]{KMS} is proved to specify its value
and certificates.

Give an action its information-set-and-action label.
A player's sequence is its ordered list of own action labels along
a history. Perfect recall gives a single preceding sequence \(s(I)\)
at all nodes in one information set \(I\). For player \(i\), the
sequence set consists of the empty sequence and each \(s(I)a\).
Its realization vector \(x_i\) obeys
\begin{equation}\label{eq:sequenceflow}
 x_i(\varnothing)=1,\qquad x_i(s)\ge0,\qquad
 \sum_{a\in A(I)}x_i(s(I)a)=x_i(s(I))\quad\text{for every }I .
\end{equation}
Write these equalities as \(E_ix_i=e_i\), where \(E_i\) has
entries in \(\{-1,0,1\}\) and \(e_i=(1,0,\ldots,0)\).
Let \(\mathcal R_i\) be this nonnegative flow polytope.

\begin{lemma}[Behavior and realization weights]\label{lem:sequence}
The nonempty compact polytope \(\mathcal R_i\) describes exactly
the own terminal realization weights of behavioral strategies,
against every opponent strategy.
\end{lemma}
\begin{proof}
A behavior assigns each sequence the product of its own probabilities,
which satisfies \eqref{eq:sequenceflow}. Conversely, from a feasible
vector set
\[
 \sigma_i(I,a)=x_i(s(I)a)/x_i(s(I))
\]
when the denominator is positive, using any distribution otherwise.
If a parent weight is zero, nonnegativity and the flow equation force
every child weight to zero. Induction on sequence length therefore
recovers every prescribed realization weight, including zero-prefix
sequences. Each coordinate lies between zero and its parent, hence
in \([0,1]\). The constraints are closed and admit uniform behavior,
proving compactness and nonemptiness.
For any opponent behavior, a terminal probability depends on player
\(i\) only through its final own sequence product.
Thus the equivalence is opponent-independent.
\end{proof}

Preselected random pure plans are covered as well: their probabilities
of agreeing with each sequence satisfy the same flow equations, and
the behavioral reconstruction has the same terminal law against every
opponent. Exponential normal-form enumeration is unnecessary.

For a policy terminal \(z\), let \(s_i(z)\) be player \(i\)'s final
sequence and let \(k_z(\alpha)\) be the product of Nature probabilities
on its complete path. Define
\begin{equation}\label{eq:sequencematrix}
 G_{s,t}(\gamma)=
    \sum_{\substack{z:s_1(z)=s\\s_2(z)=t}}
                   k_z(\alpha)u_{1,z}(\beta).
\end{equation}
For \(x\in\mathcal R_1,y\in\mathcal R_2\), expected utility is
\begin{equation}\label{eq:sequencepayoff}
 x^{\mathsf T}G(\gamma)y .
\end{equation}
Indeed each terminal probability is
\(k_z(\alpha)x(s_1(z))y(s_2(z))\); summing its utility proves
the identity. There are \(1+\sum_{I\text{ of }i}|A(I)|\) variables
and \(1+\#\{I\text{ of }i\}\) rows for player \(i\), linear in the
explicit tree size. Each leaf contributes to one matrix location.
Leaves sharing a sequence pair are summed, so a sparse matrix has
at most the number of leaves nonzero locations.

For fixed \(x\), player 2's minimization is
\[
 \min_{y\ge0,\ E_2y=e_2}(G^{\mathsf T}x)^{\mathsf T}y .
\]
It is feasible and bounded. LP strong duality identifies its value
with \(\max_v e_2^{\mathsf T}v\) subject to
\(E_2^{\mathsf T}v\le G^{\mathsf T}x\), where \(v\) is free.
Maximizing also over \(x\) gives the primal-dual pair
\begin{equation}\label{eq:sequenceLP}
 \begin{aligned}
 (\mathrm P)\quad&\max_{x,v} e_2^{\mathsf T}v\\
 &\text{subject to }E_1x=e_1,\ x\ge0,\quad
        E_2^{\mathsf T}v\le G^{\mathsf T}x,\ v\text{ free};\\[3pt]
 (\mathrm D)\quad&\min_{y,w} e_1^{\mathsf T}w\\
 &\text{subject to }E_2y=e_2,\ y\ge0,\quad
        E_1^{\mathsf T}w\ge Gy,\ w\text{ free}.
 \end{aligned}
\end{equation}
The signs can be checked directly: every feasible pair satisfies
\begin{equation}\label{eq:weakdual}
 e_2^{\mathsf T}v
   =y^{\mathsf T}E_2^{\mathsf T}v
   \le x^{\mathsf T}Gy
   \le x^{\mathsf T}E_1^{\mathsf T}w
   =e_1^{\mathsf T}w .
\end{equation}
Equality constraints give free dual variables, and nonnegative
realization variables give the displayed inequality directions.
The direct dual construction of \((\mathrm P)\) gives
\((\mathrm D)\).

Both LPs are feasible: choose any \(x\), respectively \(y\), and use
a dual optimizer for the bounded opponent best-response LP.
The primal objective is bounded above and the dual objective below
by the finite expected-utility range.
Strong duality yields equal attained optima \(v_\ell(\gamma)\),
and proves the minimax identity
\begin{equation}\label{eq:minimax}
 v_\ell(\gamma)
 =\max_{x\in\mathcal R_1}\min_{y\in\mathcal R_2}x^{\mathsf T}Gy
 =\min_{y\in\mathcal R_2}\max_{x\in\mathcal R_1}x^{\mathsf T}Gy .
\end{equation}
At equality in \eqref{eq:weakdual}, \(x\) guarantees at least \(v_\ell\)
against every \(y'\), and \(y\) at most \(v_\ell\) against every \(x'\).
They form a saddle pair. Conversely every root Nash payoff equals
this value.

Exact rational vectors satisfying \eqref{eq:sequenceLP} certify
\begin{equation}\label{eq:valuecertificate}
 \underline v_\ell=e_2^{\mathsf T}v
       \le v_\ell(\gamma)\le e_1^{\mathsf T}w=\overline v_\ell .
\end{equation}
Checking rows, signs, and inner products gives a finite exact
certificate. A gap at most \(\varepsilon_\ell\) certifies that
numerical error for the \emph{value}; zero gap gives its exact value.

For rational \(\gamma\), the matrix entries are computed from path
products and leaf sums with polynomial bit length and time in the
explicit tree and input. Rational LP has polynomial-bit optimal
primal and dual certificates and a polynomial-bit algorithm.
Thus the exact zero-sum value has polynomial explicit-tree computation,
without fixing structural dimension. This is not a strongly polynomial
LP assertion. The identities hold for real \(\gamma\), but unencoded
real numerical input has no asserted bit algorithm.

An arbitrary LP saddle behavior, especially behavior freely completed
at zero-realization information sets, is \emph{not} a certified
sequential assessment off path. The LP neither constructs consistent
beliefs nor imposes off-path sequential rationality.

\begin{lemma}[The LP computes the value of every SE]\label{lem:sevalue}
Every sequential equilibrium of a zero-sum policy game has expected
player-1 utility \(v_\ell(\gamma)\), and at least one exists.
Therefore, when \(W_\ell\) is that utility,
\begin{equation}\label{eq:valueimage}
 \mathcal W_I(\ell;p)
       =\{v_\ell(\gamma):\gamma\in\Gamma_I(p)\}.
\end{equation}
\end{lemma}
\begin{proof}
A sequential equilibrium is root Nash. To prove this directly,
partition each player's initially reachable paths by its first own
information set. The probabilities of reaching those sets depend only
on Nature and opponents, not on its deviation. At a positive-reach
first set, consistency gives the actual Bayesian conditional belief,
and \eqref{eq:sr} compares the entire alternative continuation.
Multiply by those reach probabilities and add to prove the root
Nash inequality. Paths on which the player never moves have unchanged
payoff; a zero-reach first set remains unreachable under its deviation.

The root behavior gives feasible realizations by Lemma~\ref{lem:sequence}.
The saddle value \eqref{eq:minimax} forces its zero-sum utility to
equal \(v_\ell\). Existence follows from Lemma~\ref{lem:existence},
not from a false assertion about arbitrary off-path LP completion.
Taking all compatible parameters proves \eqref{eq:valueimage}.
\end{proof}

\section{Value stability without equilibrium-selector continuity}
\label{sec:valuestability}
Fix one policy tree and two parameters \(\gamma,\gamma'\in\Gamma\).
Let \(H_\ell^{\rm N}\) be the maximum number of Nature draws along
a terminal path, including initial states, signals, and transitions.
Define
\begin{equation}\label{eq:perturbdefs}
 \begin{aligned}
 \delta_c&=\max_{\text{Nature nodes }h}
       \sum_{e\text{ leaving }h}|c_e(\alpha)-c_e(\alpha')|,\\
 \delta_u&=\max_z|u_{1,z}(\beta)-u_{1,z}(\beta')|,\qquad
 U_\ell=\max_{\gamma\in\Gamma,z}|u_{1,z}(\beta)|.
 \end{aligned}
\end{equation}
The maximum over no Nature nodes is zero; rational LPs compute
\(U_\ell\) exactly.

\begin{lemma}[Uniform fixed-profile perturbation]\label{lem:coupling}
For every pair of behavioral strategies, held fixed as functions
on the original information sets, the two expected utilities differ
by at most
\begin{equation}\label{eq:profileperturb}
 \delta_u+U_\ell H_\ell^{\rm N}\delta_c .
\end{equation}
\end{lemma}
\begin{proof}
Couple the plays while their histories agree, using the same strategic
random draws at identical information sets. At each Nature node use
a maximal coupling of its two finite rows, whose failure probability
is their total variation and at most \(\delta_c/2\).
There are at most \(H_\ell^{\rm N}\) opportunities for a first failure.
Summing conditional first-failure probabilities bounds the probability
of different terminal histories by \(H_\ell^{\rm N}\delta_c/2\).
At identical terminals the utility difference is at most \(\delta_u\).
At different terminals, separate the parameter change at one terminal,
at most \(\delta_u\), from the difference between two utilities at
\(\gamma'\), at most \(2U_\ell\).
Taking expectation proves \eqref{eq:profileperturb}.
Private persistent information is preserved: before histories differ,
the same history uses the same information-set distribution.
\end{proof}

If two functions differ uniformly by at most \(\varepsilon\),
their minima and maxima differ by at most \(\varepsilon\):
the pointwise inequalities imply both one-sided bounds after
optimization. Apply this fact twice to \eqref{eq:minimax} and
Lemma~\ref{lem:coupling} to get
\begin{equation}\label{eq:valuestability}
 |v_\ell(\gamma)-v_\ell(\gamma')|
                \le\delta_u+U_\ell H_\ell^{\rm N}\delta_c .
\end{equation}
No unique equilibrium, strict incentive gap, nonsingular policy
Jacobian, or continuous selector is used.
For Nature coefficient rows \(B_e\), set
\begin{equation}\label{eq:valueLipschitzconstants}
 \begin{aligned}
 L_{c,\ell}&=\max_h\sum_{e\text{ leaving }h}\|B_e\|_1,\\
 L_{u,\ell}&=\max_z\|\varphi_{1,z}\|_1,\qquad
 L_\ell=L_{u,\ell}+U_\ell H_\ell^{\rm N}L_{c,\ell}.
 \end{aligned}
\end{equation}
These explicit constants give
\begin{equation}\label{eq:valueLipschitz}
 \begin{split}
 |v_\ell(\gamma)-v_\ell(\gamma')|
 &\le L_{u,\ell}\|\beta-\beta'\|_\infty
       +U_\ell H_\ell^{\rm N}L_{c,\ell}\|\alpha-\alpha'\|_\infty\\
 &\le L_\ell\|\gamma-\gamma'\|_\infty .
 \end{split}
\end{equation}
Repeated rows and joint parameter restrictions do not invalidate
these inequalities; sharper rowwise constants are optional.

\section{Polynomial confidence enclosures for policy values}
\label{sec:valueconfidence}
Suppose the statistical construction returns a rational
\(\mathcal C_n\subseteq\Gamma\) and a common measurable coverage
event \(E_n\) on which \(\Gamma_I(p)\subseteq\mathcal C_n\).
On that event the polytope is nonempty.
When nonempty, choose a \emph{fresh} rational
\(\gamma_{\rm center}\in\mathcal C_n\) by LP.
It is distinct from the minimax-fit \(\widehat\gamma\):
the fit uses one common residual tolerance, while
\eqref{eq:confidencepoly} uses different block tolerances, so
\(\widehat\gamma\) need not belong to \(\mathcal C_n\).

Compute a certified bound
\[
 D_n\ge\sup_{\gamma\in\mathcal C_n}
                  \|\gamma-\gamma_{\rm center}\|_\infty.
\]
The \(2d\) coordinate LPs, for ambient dimension \(d\), give such
a radius exactly by taking the largest distance from the center
to any coordinate endpoint. Any certified diameter upper bound
also suffices because the center belongs to \(\mathcal C_n\).
An empty set is reported as empty and cannot arise on the coverage
event under a correctly specified nonempty population model.

For each zero-sum policy, compute at \(\gamma_{\rm center}\)
feasible LP bounds
\([\underline v_\ell,\overline v_\ell]\) with gap at most
\(\varepsilon_\ell\). Equations \eqref{eq:valuecertificate},
\eqref{eq:valueimage}, and \eqref{eq:valueLipschitz} imply that
\begin{equation}\label{eq:valueinterval}
 I_{\ell,n}=
  [\,\underline v_\ell-L_\ell D_n,\
                 \overline v_\ell+L_\ell D_n\,]
\end{equation}
contains every value in \(\mathcal W_I(\ell;p)\) for every declared
policy on the same event \(E_n\). Its length is at most
\begin{equation}\label{eq:valueintervallength}
 \varepsilon_\ell+2L_\ell D_n .
\end{equation}
This separates LP numerical error from sampling uncertainty.
The interval is a certified enclosure, not a claim of sharpness
over \(\mathcal C_n\).

For two policy values, bracket
\(v_\ell(\gamma_{\rm center})-v_{\ell'}(\gamma_{\rm center})\),
then enlarge each side by \((L_\ell+L_{\ell'})D_n\).
The numerical gap is at most
\(\varepsilon_\ell+\varepsilon_{\ell'}\).
The same parameter is kept. Selection is immaterial only because
each target is the zero-sum utility value, not because policy
equilibria are unique.

Under the full-rank statistical condition, on
\(\Delta\le e_n\), equation \eqref{eq:S12} gives
\[
 D_n\le\diam_2(\mathcal C_n)
          \le4\sqrt R\,Ce_n/\kappa
\]
if the exact coordinate radius is used.
Indeed each coordinate distance is bounded by the true diameter.
Thus \eqref{eq:valueintervallength} is an honest root-\(n\)
enclosure at fixed input, confidence, and positive conditioning.
Choose \(\varepsilon_\ell\) of the same or smaller order.
Weak conditioning enlarges the constants explicitly.
All computations for these intervals are polynomial in the explicit
tree and rational confidence-polytope input; this makes no succinct-
horizon claim.

\section{A limit to general outcome contraction at every confidence level}
\label{sec:impossibility}
Parameter rates alone do not imply continuity of a policy's entire
equilibrium-outcome range. The following example lies inside the
declared affine class.
Let \(\theta\in[-1/4,1/4]\) and let baseline Nature draw
\(Y\sim\operatorname{Bernoulli}(1/2+\theta)\).
Baseline utilities are zero. If an action stage is desired, one active
player chooses \(a\in\{0,1\}\), still with zero utility, and the same
independent fair mixed behavior is selected at every parameter
and across markets. The researcher observes \((Y,a)\) completely.
The ancillary action supplies no information about \(\theta\).
The identity measurement is full rank, all complete-history probabilities
are at least \(1/8\), and Nature probabilities are uniformly positive.
The identified parameter is
\[
 \theta=\Prob(Y=1)-1/2,\qquad
 \widehat\theta=\overline Y-1/2,\qquad
 \E(\widehat\theta-\theta)^2\le1/(4n).
\]
Ordinary Hoeffding intervals give uniform root-\(n\) coverage.

The policy changes only the active player's utility to \(\theta a\),
and the outcome statistic is \(W=a\). Its complete range over
sequential equilibria is
\begin{equation}\label{eq:tierange}
 \mathcal R(\theta)=
 \begin{cases}
  \{0\},&\theta<0,\\
  [0,1],&\theta=0,\\
  \{1\},&\theta>0.
 \end{cases}
\end{equation}
Strict preference forces the nonzero-sign cases at every information
set. At zero, every behavior is rational and admits fixed-game
completely mixed approximations with the corresponding Bayes beliefs;
every expected action between zero and one is attained.

Use the joint rational polytope
\(\alpha=\beta=\theta\in[-1/4,1/4]\).
Its equality transfers the identified chance coefficient to the
payoff coefficient, so the example meets the common model even though
an ambient-rank test need not be necessary. A dummy second player
with utility \(-\theta a\) embeds it into a zero-sum game.
Its utility value \(\max\{0,\theta\}\) is Lipschitz, while the
different outcome \(a\) has the discontinuous range
\eqref{eq:tierange}. Zero-sum structure does not make every outcome
equal to the value.

\begin{theorem}[No uniformly contracting honest general-outcome sets]
\label{thm:impossibility}
Fix any \(0<\eta<1\). Let \(\mathcal D_n\) be measurable random
closed sets with whole-range finite-sample honesty
\begin{equation}\label{eq:honesttie}
 \inf_{\theta\in[-1/4,1/4]}
   \Prob_\theta\{\mathcal R(\theta)\subseteq\mathcal D_n\}
                         \ge1-\eta.
\end{equation}
For \(\theta_n=1/n\), \(n\ge4\),
\begin{equation}\label{eq:tielower}
 \begin{split}
 \Prob_{\theta_n}\{\diam(\mathcal D_n)\ge1,\
     \excess(\mathcal D_n,\mathcal R(\theta_n))\ge1\}
 \ \ge\ 1-\eta-\tfrac12\sqrt{(1+4/n^2)^n-1}.
 \end{split}
\end{equation}
Here
\(\excess(A,B)=\sup_{x\in A}\inf_{y\in B}|x-y|\).
Thus Hausdorff error cannot vanish uniformly in probability.
The same conclusion holds under asymptotic uniform honesty, with
an additional vanishing slack.
\end{theorem}
\begin{proof}
Let \(P_0,P_\theta\) be the one-market Bernoulli laws.
Their likelihood ratio under \(P_0\) takes values
\(1+2\theta\) and \(1-2\theta\). Hence
\[
 \E_0\left[\left(\frac{dP_\theta}{dP_0}\right)^2\right]
      =1+4\theta^2,\qquad
 \chi^2(P_\theta^n,P_0^n)=(1+4\theta^2)^n-1.
\]
The product identity follows by independence. Cauchy--Schwarz
applied to the centered likelihood ratio gives
\[
 \TV(P_\theta^n,P_0^n)
          \le\tfrac12\sqrt{(1+4\theta^2)^n-1}.
\]
An independent ancillary action leaves this distance unchanged.
At \(\theta_n=1/n\) the bound tends to zero.

At zero, honesty says the single event
\(A_n=\{[0,1]\subseteq\mathcal D_n\}\) has probability
at least \(1-\eta\).
The probabilities of that same event under \(P_0^n\) and
\(P_{\theta_n}^n\) differ by at most their total variation,
giving the right-hand side of \eqref{eq:tielower}.
On \(A_n\), the reported set contains zero and one, while
\(\mathcal R(\theta_n)=\{1\}\). Its diameter and its excess
over the true range are both at least one, proving the result.
The excess also lower-bounds the Hausdorff error.
Asymptotic uniform honesty repeats the argument with its vanishing
coverage slack. There is no intersection of two confidence events,
so the proof works for every \(0<\eta<1\).
\end{proof}

The obstruction is equilibrium indifference, not a discontinuous
primitive statistic or vanishing Nature probability. The terminal
statistic \(a\) is bounded and polynomial. General policy outcomes
therefore retain the exact range and confidence-coverage methods;
they do not inherit the value contraction theorem.

\section{Scope, attribution, and a rational illustration}
\label{sec:conclusion}
The results retain every baseline and policy sequential assessment
in the declared fixed-tree, positive-Nature class.
The generic exact branch includes rank-deficient observation channels
and strategically zero baseline histories.
The informative branch provides a proved observation restriction,
compatibility test, and sharp affine slice for arbitrary finite known
private-history trees. Its complete proofs, not just the Bayesian
family, yield the inference results.

The statistical construction separates three questions.
The minimax parameter fit estimates admitted primitives.
The outer polytope covers the full sharp parameter fiber.
The reconstructed profile has exact feasibility and Bayes consistency,
with a computed certificate for all continuation deviations.
The latter is not a replacement for exact policy-equilibrium coverage.
General counterfactual sets retain exact endpoints and numerical
error certificates, while the fast contraction result concerns
zero-sum utility values. Comparative effects keep a shared structural
parameter and either all policy pairs or an explicitly supplied
joint selection relation.

All polynomial-time claims specify an explicit rational input.
The complete history tree can be exponentially larger than a compact
state/horizon description. Full-history measurement calibration is
substantive. Unknown labels, hidden equilibrium mixtures, missing
payoff normalization, weak support, or weak rank cannot be removed by
reporting a successful optimizer. The bounds display where those
conditions enter. Varying primitive support, arbitrary unencoded
utilities, and infinite horizons require different theorems.

As one rational instance of Proposition~\ref{prop:bayesian}, take
\[
 \mathsf Q=
   \begin{pmatrix}7/16&1/16\\1/16&7/16\end{pmatrix},\qquad
 \theta=(3/4,5/4),\qquad \xi=1/2.
\]
Then
\[
 p^a=(4/7,4/9),\qquad q^a=(16/27,80/189),\qquad
 v=31/63 .
\]
There are sixteen complete terminal histories, and their minimum
mass is \(11/972\ge1/144=q_{\min}/9\) with \(q_{\min}=1/16\).
The inverse channel and conditional equations recover the same
labeled type law and payoff coefficients.
These exact arithmetic identities illustrate the general proof;
they do not constitute a full run of the QE, LP, or statistical
pipeline, and no such execution is claimed here.

The exact SE/QE baseline, compactness, endpoint method, and elementary
whole-set coverage are retained from the prior attempt \cite{Prior}.
The general informative fiber and its conditioning, diameter, and
actual-assessment certificates supply the additional obligations
left there. Sequence form, real quantifier elimination, Brouwer's
fixed-point theorem, and rational LP strong duality and polynomial-bit
algorithms are classical inputs. The manuscript supplies the
reductions, equivalences, and bounds needed for their present use.
It makes no claim of external peer review, formal proof-assistant
verification, or historical novelty.

\begin{thebibliography}{9}
\bibitem{KW}
D.~M. Kreps and R. Wilson,
\emph{Sequential Equilibria},
Econometrica 50 (1982), 863--894.
\url{https://doi.org/10.2307/1912767}.
Institutional bibliographic record:
\url{https://www.gsb.stanford.edu/faculty-research/publications/sequential-equilibrium}.
Used for the equilibrium concept; the fixed-game encoding and
equivalences used here are proved in Section~\ref{sec:se}.

\bibitem{Basu}
S. Basu, \emph{Algorithms in Real Algebraic Geometry: A Survey},
arXiv:1409.1534v1 (4 September 2014).
\url{https://arxiv.org/pdf/1409.1534v1}.
Inspected locators: Theorem 2.27, printed p.~16, including its
integer bit-size clause; Theorem 2.4, printed p.~6.
These supply the imported effective quantifier-elimination and
fixed-compiled-size complexity bounds.

\bibitem{KMS}
D. Koller, N. Megiddo, and B. von Stengel,
\emph{Efficient Computation of Equilibria for Extensive Two-Person Games},
Games and Economic Behavior 14 (1996), 247--259.
\url{https://ai.stanford.edu/~koller/Papers/Koller%2Bal%3AGEB96.pdf}.
Inspected locator: Section 2, pp.~249--250, for realization flows,
behavioral reconstruction, Nature-weighted terminal matrices, and
best-response primal/dual LPs. No polynomial-time general-sum LCP
claim is imported.

\bibitem{Prior}
GPT 6 Astra Ultra,
\emph{Exact finite-game identification with consistent off-path beliefs},
prior catalogue attempt C73-49, 9 October 2026,
at pinned repository commit
\texttt{00ced55ba1223fb48603e1f15028d753997ee515}.
\href{https://github.com/mehmetmars7/Erdosproblems-llm-hunter/blob/00ced55ba1223fb48603e1f15028d753997ee515/attacks/open_problems/economics/gpt6_astra_ultra/C73-49.tex}
{Pinned complete prior source}.
Its exact SE reduction, compactness, sharp policy projections,
endpoint bisection, and whole-set coverage are reproduced with
attribution. Informative reconstruction and quantitative statistical
and value certificates are the additional obligations addressed here.
\end{thebibliography}
\end{document}
